\chapter{Particles, Fields, and Relativistic Notation}

\section{A Doppler-Shift Problem Before the Lecture}

A question left over from the preceding week opens the discussion.  Does an
observer moving nearly at the speed of light always measure a longer
wavelength?  No: the answer depends on the direction of motion.  An observer
moving with a light wave sees its wavelength stretched, whereas an observer
moving against the wave sees it compressed.  This is the relativistic Doppler
effect.

\begin{classroomqa}
\audiencequestion{Does motion near the speed of light make the observed wavelength longer?}
\lecturerresponse{Only when the observer moves in the same direction as the light.  Reversing the observer's velocity reverses the shift.  The result can be obtained directly from Lorentz geometry, without memorizing a separate formula. \lecturetimestamp{00:00:13}}
\end{classroomqa}

Use units with \(c=1\), and let the primed frame move to the right with speed
\(v\).  A right-moving light crest follows a null line.  Normalize the
unprimed wavelength to one and write two successive crest worldlines as
\begin{equation}
x=t,
\qquad
x=t+1.
\label{eq:sr5-crest-lines}
\end{equation}
The wavelength in the unprimed frame is measured between the crests at equal
\(t\).  The wavelength in the primed frame must instead be measured between
them at equal \(t'\).  Relativity of simultaneity is therefore the whole
calculation.

\begin{figure}[htbp]
\centering
\includegraphics[width=0.90\textwidth]{lecture_05_figure_01.png}
\caption{Successive right-moving wave crests, the light cone, and the tilted axes of the moving frame.}
\lectureframe{00:03:30}
\end{figure}

At this point the derivation pauses for a short discussion of the course
sequence.  There will be no summer course.  General relativity and a broader
return to quantum mechanics are both candidates for the next course; the
harmonic oscillator is singled out as material that should be understood
before quantum field theory.

\begin{classroomqa}
\audiencequestion{Could a future course discuss the infrared--ultraviolet connection?}
\lecturerresponse{That question belongs to the meeting of quantum mechanics and gravity.  The class has enough quantum mechanics to recognize the issue, but it must first develop gravity, most naturally through general relativity. \lecturetimestamp{00:05:37}}
\end{classroomqa}

\begin{classroomqa}
\audiencequestion{Can the relative scale of the tilted axes simply be read from the blackboard drawing?}
\lecturerresponse{No.  A spacetime diagram records causal and simultaneity geometry, but an ordinary Euclidean ruler on the drawing is not a Lorentz measurement.  Locate the relevant events and apply the Lorentz transformation. \lecturetimestamp{00:06:10}}
\end{classroomqa}

The calculation now resumes.  Normalize the unprimed crest spacing to one and
locate the two simultaneous primed-frame events algebraically.

The Lorentz transformation is
\begin{equation}
t'=\gamma(t-vx),
\qquad
x'=\gamma(x-vt),
\qquad
\gamma=\frac{1}{\sqrt{1-v^2}}.
\label{eq:sr5-lorentz-doppler}
\end{equation}
Choose the primed simultaneity slice \(t'=0\).  Equation
\eqref{eq:sr5-lorentz-doppler} then gives
\begin{equation}
t=vx.
\label{eq:sr5-primed-simultaneity}
\end{equation}
The first crest intersects this slice at the origin.  The second crest obeys
\(x=t+1\), so its intersection satisfies
\begin{align}
x&=vx+1,\\
x&=\frac{1}{1-v},
\qquad
t=\frac{v}{1-v}.
\end{align}
Its primed coordinate is
\begin{align}
x'
&=\frac{x-vt}{\sqrt{1-v^2}}\\
&=\frac{1-v^2}{(1-v)\sqrt{1-v^2}}\\
&=\sqrt{\frac{1+v}{1-v}}.
\end{align}
Since the original wavelength was one,
\begin{equation}
\boxed{
\frac{\lambda'}{\lambda}
=
\sqrt{\frac{1+\beta}{1-\beta}}
},
\qquad
\beta=\frac{v}{c}.
\label{eq:sr5-doppler}
\end{equation}
Equivalently,
\begin{equation}
\frac{\lambda'}{\lambda}
=
\sqrt{\frac{c+v}{c-v}}.
\end{equation}

\begin{figure}[htbp]
\centering
\includegraphics[width=0.90\textwidth]{lecture_05_figure_02.png}
\caption{The Doppler factor after restoring \(c\); reversing \(v\) gives the reciprocal shift.}
\lectureframe{00:13:40}
\end{figure}

For \(v>0\), the primed observer chases the light and measures a longer
wavelength.  Replacing \(v\) by \(-v\) describes motion against the light and
makes the factor smaller than one.  If an observer could approach \(c\) while
chasing the ray, the wavelength would diverge.  The frequency shift is the
inverse ratio.

\section{Why Interaction Must Appear in One Action}

After the preliminary questions, the subject is fields and particles.  In
quantum field theory, particles and fields
will eventually become two descriptions of the same excitations.  That is not
the framework here.  Here particles follow worldlines and classical fields
fill spacetime.  The immediate question is more modest: if a field acts on a
particle, how does the particle act back on the field?

The action principle answers this through a finite-dimensional prototype.
Let \(x(t)\) and \(y(t)\) be two generalized coordinates.  If the Lagrangian
separates,
\begin{equation}
L(x,\dot x,y,\dot y)
=
L_x(x,\dot x)+L_y(y,\dot y),
\label{eq:sr5-separable-l}
\end{equation}
then the two Euler--Lagrange equations are independent:
\begin{equation}
\frac{d}{dt}\frac{\partial L}{\partial\dot x}
=
\frac{\partial L}{\partial x},
\qquad
\frac{d}{dt}\frac{\partial L}{\partial\dot y}
=
\frac{\partial L}{\partial y}.
\end{equation}
The \(x\)-equation contains no \(y\), and the \(y\)-equation contains no \(x\).

To make the variables interact, add a term that cannot be split into an
\(x\)-part and a \(y\)-part.  A simple example is
\begin{equation}
L
=
\frac{\dot x^2}{2}
+\frac{\dot y^2}{2}
-V(x)-W(y)+\kappa xy.
\label{eq:sr5-coupled-coordinates}
\end{equation}
Now varying \(x\) produces a term proportional to \(y\), while varying \(y\)
produces a term proportional to \(x\).  The same mixed term governs both
directions of influence.

\begin{figure}[htbp]
\centering
\includegraphics[width=0.90\textwidth]{lecture_05_figure_03.png}
\caption{Two generalized coordinates become dynamically coupled when one term in the Lagrangian contains both of them.}
\lectureframe{00:19:00}
\end{figure}

This is the precise setting of the action--reaction statement used here: both
objects are dynamical and both are varied in one common action.  Treating a
field as an externally prescribed background can hide its response, but once
the field itself is dynamical the coupling term supplies both equations.

\section{How a Scalar Field Acts on a Particle}

\subsection{The relativistic coupling}

Begin by treating the scalar field \(\phi(\mathbf{x},t)\) as known.  It may be
a wave or a static profile; for the moment its own equation is irrelevant.
The free relativistic particle action is
\begin{equation}
S_{\mathrm p}
=
-m\int d\tau.
\end{equation}
For motion on a line and \(c=1\),
\begin{equation}
L_{\mathrm p}^{(0)}
=
-m\sqrt{1-\dot x^2}.
\end{equation}
The simplest scalar coupling replaces the mass coefficient by a field-dependent
one:
\begin{equation}
\boxed{
L_{\mathrm p}
=
-\bigl[m+g\phi(x,t)\bigr]\sqrt{1-\dot x^2}
}.
\label{eq:sr5-particle-coupling}
\end{equation}
The coupling constant \(g\) measures how strongly the field affects this
particle.  Its sign is a convention: changing \(g\) to \(-g\) reverses the
direction of the scalar force.

\begin{figure}[htbp]
\centering
\includegraphics[width=0.90\textwidth]{lecture_05_figure_04.png}
\caption{The free particle action with \(m\) replaced by the local coefficient \(m+g\phi(x,t)\).}
\lectureframe{00:24:10}
\end{figure}

\subsection{The slow-particle limit}

Restore the speed of light so that the Newtonian limit is visible:
\begin{equation}
L_{\mathrm p}
=
-\bigl[mc^2+g\phi(x,t)\bigr]
\sqrt{1-\frac{\dot x^2}{c^2}}.
\label{eq:sr5-particle-coupling-c}
\end{equation}
For \(|\dot x|\ll c\),
\begin{equation}
\sqrt{1-\frac{\dot x^2}{c^2}}
=
1-\frac{\dot x^2}{2c^2}
+O\!\left(\frac{\dot x^4}{c^4}\right).
\end{equation}
Substitution into \eqref{eq:sr5-particle-coupling-c} gives
\begin{align}
L_{\mathrm p}
&=
-mc^2-g\phi
+\frac{m\dot x^2}{2}
+\frac{g\phi\,\dot x^2}{2c^2}
+O(c^{-4})\\
&\longrightarrow
\frac{m\dot x^2}{2}-g\phi(x,t),
\label{eq:sr5-nr-particle}
\end{align}
after dropping the constant \(-mc^2\) and terms suppressed by \(c^{-2}\).
The result is the ordinary form \(T-V\), so
\begin{equation}
V_{\mathrm p}(x,t)=g\phi(x,t).
\label{eq:sr5-scalar-potential}
\end{equation}

\begin{figure}[htbp]
\centering
\includegraphics[width=0.90\textwidth]{lecture_05_figure_05.png}
\caption{The square-root expansion leaves the Newtonian kinetic term and the scalar potential energy \(g\phi\).}
\lectureframe{00:28:50}
\end{figure}

The large rest-energy constant does not enter the Euler--Lagrange equation.
The field-dependent term does.  In electromagnetism, electric charge plays
the analogous role of a coupling strength, although the actual
electromagnetic interaction uses a four-vector potential rather than this
scalar toy model.

\section{How the Particle Acts Back on the Field}

\subsection{One variational problem}

The field cannot now remain merely prescribed.  Consider a spacetime region
containing both a field configuration and a particle worldline.  The physical
solution makes the total action stationary under variations of both:
\begin{equation}
\delta S[\phi,x]=0.
\end{equation}
Varying the worldline while holding the field fixed gives the particle
equation.  Varying the field while holding the worldline fixed gives the field
equation.  Their common interaction term is what joins them.

\begin{figure}[htbp]
\centering
\includegraphics[width=0.90\textwidth]{lecture_05_figure_06.png}
\caption{A particle worldline passing through a spacetime region occupied by a scalar field.}
\lectureframe{00:30:30}
\end{figure}

For the free massless scalar field, take
\begin{equation}
S_{\mathrm f}
=
\int d^4x\,
\frac12
\left[
(\partial_t\phi)^2
-(\partial_x\phi)^2
-(\partial_y\phi)^2
-(\partial_z\phi)^2
\right].
\label{eq:sr5-free-field-action}
\end{equation}
From this point onward, use \(c=1\) again.  With explicit units, the time term would
carry the corresponding factor of \(1/c^2\), depending on the normalization
of the coordinates and field.

\begin{figure}[htbp]
\centering
\includegraphics[width=0.90\textwidth]{lecture_05_figure_07.png}
\caption{The field action written over the same spacetime region as the particle worldline.}
\lectureframe{00:32:10}
\end{figure}

\subsection{A static particle and a worldline integral}

To expose the source with minimal algebra, place the particle at rest at the
origin.  The factor \(d\tau\) then equals \(dt\), and the interaction part of
\eqref{eq:sr5-particle-coupling} becomes
\begin{equation}
S_{\mathrm{int}}
=
-g\int dt\,\phi(\mathbf{0},t).
\label{eq:sr5-worldline-interaction}
\end{equation}
This is an integral only along the worldline, whereas
\eqref{eq:sr5-free-field-action} is an integral over all spacetime.  To combine
them into one spacetime integral, we need a function that selects the
particle's position.

\begin{figure}[htbp]
\centering
\includegraphics[width=0.90\textwidth]{lecture_05_figure_08.png}
\caption{The static worldline inside the field region and the field action that must be combined with it.}
\lectureframe{00:35:35}
\end{figure}

\subsection{The Dirac delta function}

The one-dimensional delta function is defined by
\begin{equation}
\int_{-\infty}^{\infty}dx\,f(x)\delta(x-a)=f(a).
\label{eq:sr5-delta-definition}
\end{equation}
It vanishes away from \(x=a\), is singular at \(x=a\), and has unit integral.
It is best understood as a distribution characterized by
\eqref{eq:sr5-delta-definition}, rather than as an ordinary finite-valued
function.

\begin{figure}[htbp]
\centering
\includegraphics[width=0.90\textwidth]{lecture_05_figure_09.png}
\caption{The delta-function identity that rewrites the field value at the origin as a spatial integral.}
\lectureframe{00:36:45}
\end{figure}

In three dimensions,
\begin{equation}
\delta^3(\mathbf{x})
=
\delta(x)\delta(y)\delta(z),
\end{equation}
and therefore
\begin{align}
\int d^3x\,
\phi(\mathbf{x},t)\delta^3(\mathbf{x})
&=
\int dx\,dy\,dz\,
\phi(x,y,z,t)\\
&\quad\times\delta(x)\delta(y)\delta(z)\\
&=\phi(\mathbf{0},t).
\label{eq:sr5-triple-delta}
\end{align}

\begin{classroomqa}
\audiencequestion{Where does the \(d^3x\) come from?}
\lecturerresponse{It is the ordinary spatial volume element \(dx\,dy\,dz\).  In one spatial dimension the same identity contains only \(dx\) and one delta function. \lecturetimestamp{00:37:09}}
\end{classroomqa}

\begin{figure}[htbp]
\centering
\includegraphics[width=0.90\textwidth]{lecture_05_figure_10.png}
\caption{The three-dimensional delta function expanded as three one-dimensional factors.}
\lectureframe{00:39:05}
\end{figure}

Equation \eqref{eq:sr5-worldline-interaction} can now be written as
\begin{equation}
S_{\mathrm{int}}
=
-g\int d^4x\,
\delta^3(\mathbf{x})\phi(\mathbf{x},t).
\end{equation}
The complete static-source action is
\begin{align}
S
=
\int d^4x
\biggl\{
&\frac12(\partial_t\phi)^2
-\frac12|\boldsymbol{\nabla}\phi|^2
\nonumber\\
&-g\delta^3(\mathbf{x})\phi
\biggr\}.
\label{eq:sr5-combined-action}
\end{align}

\begin{classroomqa}
\audiencequestion{Was a minus sign lost when the interaction was inserted into the spacetime action?}
\lecturerresponse{The worldline Lagrangian contains \(-g\phi\), so the spacetime density must contain \(-g\delta^3(\mathbf{x})\phi\).  Keeping that sign fixes the source term in the field equation. \lecturetimestamp{00:41:24}}
\end{classroomqa}

\begin{figure}[htbp]
\centering
\includegraphics[width=0.90\textwidth]{lecture_05_figure_11.png}
\caption{The free scalar density and the localized interaction assembled into one spacetime action.}
\lectureframe{00:42:05}
\end{figure}

\subsection{The sourced field equation}

Varying \eqref{eq:sr5-combined-action} with respect to \(\phi\) gives
\begin{equation}
\boxed{
\partial_t^2\phi-\nabla^2\phi
=
-g\delta^3(\mathbf{x})
}.
\label{eq:sr5-sourced-wave}
\end{equation}
Away from the origin this is the free wave equation.  At the origin the delta
function marks the particle as a localized source.

\begin{figure}[htbp]
\centering
\includegraphics[width=0.90\textwidth]{lecture_05_figure_12.png}
\caption{The wave operator driven by a point source at the origin.}
\lectureframe{00:45:35}
\end{figure}

For a static solution, \(\partial_t\phi=0\), so
\begin{equation}
\boxed{
\nabla^2\phi
=
g\delta^3(\mathbf{x})
}.
\label{eq:sr5-poisson}
\end{equation}
This is Poisson's equation for a point source.  The analogy is electrostatic,
not yet electromagnetic: the field is still the scalar \(\phi\), not the
electromagnetic four-potential.

\begin{figure}[htbp]
\centering
\includegraphics[width=0.90\textwidth]{lecture_05_figure_13.png}
\caption{The static limit of the sourced wave equation, recognized as Poisson's equation.}
\lectureframe{00:47:05}
\end{figure}

The reciprocity is now explicit.  The term \(g\phi\) gives the particle a
potential energy, and the same \(g\) makes the particle a source for \(\phi\).
The two effects are not independent assumptions.

\subsection{Moving sources radiate}

A particle fixed at \(\mathbf{a}\) is represented by
\(\delta^3(\mathbf{x}-\mathbf{a})\).  If its position changes with time, the
source becomes
\begin{equation}
\delta^3\!\left(\mathbf{x}-\mathbf{a}(t)\right).
\label{eq:sr5-moving-delta}
\end{equation}
A time-dependent right-hand side cannot be balanced by a time-independent
field.  The \(\partial_t^2\phi\) term must return, and an accelerating or
oscillating source can launch waves into the field.

\begin{figure}[htbp]
\centering
\includegraphics[width=0.90\textwidth]{lecture_05_figure_14.png}
\caption{A delta-function source shifted to \(\mathbf{a}(t)\), making the field equation time dependent.}
\lectureframe{00:50:30}
\end{figure}

\begin{classroomqa}
\audiencequestion{If the source is treated nonrelativistically, where is the relativity in the field equation?}
\lecturerresponse{The relativistic propagation is encoded in the spacetime wave operator on the left.  A fully relativistic particle treatment would retain its proper-time action and velocity factors.  Even a slowly moving but time-dependent source can drive waves. \lecturetimestamp{00:52:14}}
\end{classroomqa}

\begin{classroomqa}
\audiencequestion{What units must \(g\phi\) have?}
\lecturerresponse{It appears beside \(mc^2\) in the particle action, so it has units of energy.  In the natural-unit normalization used here, a scalar in \(3+1\) dimensions has mass dimension one and \(g\) may be dimensionless.  The dimensions change with spacetime dimension and with field normalization. \lecturetimestamp{00:52:55}}
\end{classroomqa}

\section{Four-Vectors, Scalars, and the Metric}

The component equations have become long enough to justify a better notation.
Write the spacetime coordinates as
\begin{equation}
x^\mu=(x^0,x^1,x^2,x^3)=(t,x,y,z),
\label{eq:sr5-four-coordinate}
\end{equation}
where Greek indices run over all four components.  A four-component list is
not automatically a four-vector.  It is a four-vector only if its components
transform under Lorentz transformations as \(x^\mu\) does.

\begin{figure}[htbp]
\centering
\includegraphics[width=0.90\textwidth]{lecture_05_figure_15.png}
\caption{The four spacetime coordinates collected into \(x^\mu\).}
\lectureframe{00:55:00}
\end{figure}

A statement that the three spatial components vanish may hold in one frame
and fail in another.  By contrast, if all four components vanish, every
Lorentz observer agrees that the vector is zero.  Differences such as
\(dx^\mu\) are also four-vectors.

The invariant proper-time interval is
\begin{equation}
d\tau^2
=
dt^2-dx^2-dy^2-dz^2.
\label{eq:sr5-proper-time}
\end{equation}
It has one value, not four frame-dependent components, so it is a scalar.

\begin{figure}[htbp]
\centering
\includegraphics[width=0.90\textwidth]{lecture_05_figure_16.png}
\caption{The four-displacement and the invariant proper-time interval.}
\lectureframe{00:57:25}
\end{figure}

\subsection{The Minkowski metric}

Adopt the mostly-plus signature
\begin{equation}
\eta_{\mu\nu}
=
\operatorname{diag}(-1,1,1,1).
\label{eq:sr5-metric}
\end{equation}
The alternative mostly-minus convention is equally valid, but every equation
must use one convention consistently.  With \eqref{eq:sr5-metric},
\begin{equation}
dx^\mu dx_\mu
=
-dt^2+dx^2+dy^2+dz^2
=
-d\tau^2.
\end{equation}
The overall sign does not alter the invariance.

\begin{figure}[htbp]
\centering
\includegraphics[width=0.90\textwidth]{lecture_05_figure_17.png}
\caption{The mostly-plus Minkowski metric written as a diagonal matrix.}
\lectureframe{01:00:30}
\end{figure}

The metric lowers an index:
\begin{equation}
A_\mu=\eta_{\mu\nu}A^\nu.
\label{eq:sr5-lower-index}
\end{equation}
For this Cartesian metric,
\begin{equation}
A^\mu=(A^0,A^1,A^2,A^3),
\qquad
A_\mu=(-A^0,A^1,A^2,A^3).
\label{eq:sr5-upper-lower-components}
\end{equation}
Applying the same metric again raises the index because
\(\eta^{\mu\rho}\eta_{\rho\nu}=\delta^\mu{}_\nu\).

\begin{figure}[htbp]
\centering
\includegraphics[width=0.90\textwidth]{lecture_05_figure_18.png}
\caption{Index lowering as matrix multiplication by \(\eta_{\mu\nu}\).}
\lectureframe{01:03:00}
\end{figure}

\begin{classroomqa}
\audiencequestion{Does ordinary matrix multiplication really contain a summation here?}
\lecturerresponse{Yes.  Each output component is a row of \(\eta\) multiplied by the full vector: a sum over the column index.  The diagonal metric makes the calculation look trivial, but a general matrix makes the sum explicit. \lecturetimestamp{01:04:56}}
\end{classroomqa}

\subsection{Einstein summation and free indices}

Einstein's convention suppresses the summation symbol whenever an index is
repeated once upstairs and once downstairs:
\begin{equation}
\eta_{\mu\nu}A^\nu
\equiv
\sum_{\nu=0}^{3}\eta_{\mu\nu}A^\nu.
\end{equation}
Here \(\nu\) is a dummy index; it is summed away and may be renamed.  The index
\(\mu\) is free and labels the four components of the result.  An index should
not normally appear more than twice in one monomial.

The traditional names are \emph{contravariant} for an upper index and
\emph{covariant} for a lower index.  The lecture deliberately prefers the
plainer language \emph{upper} and \emph{lower}, but both terminologies refer to the
same transformation rules.

\section{Contractions and Derivatives of a Scalar}

\subsection{The Minkowski inner product}

Contracting a vector with its lowered form gives
\begin{equation}
A^\mu A_\mu
=
\eta_{\mu\nu}A^\mu A^\nu
=
-(A^0)^2+(A^1)^2+(A^2)^2+(A^3)^2.
\label{eq:sr5-vector-square}
\end{equation}
No extra minus sign should be inserted by hand; it already comes from the
lowered time component.

\begin{figure}[htbp]
\centering
\includegraphics[width=0.90\textwidth]{lecture_05_figure_19.png}
\caption{The contraction \(A^\mu A_\mu\) expanded into temporal and spatial components.}
\lectureframe{01:09:30}
\end{figure}

\begin{classroomqa}
\audiencequestion{Should a lower-index vector be drawn as a row rather than a column?}
\lecturerresponse{It may be represented as a row in elementary matrix notation, but index notation does not require that choice.  Its defining feature is how it transforms and how it contracts with an upper-index vector. \lecturetimestamp{01:10:26}}
\end{classroomqa}

\begin{classroomqa}
\audiencequestion{Do Greek indices always refer to all four components?}
\lecturerresponse{In this notation, yes.  Greek indices such as \(\mu,\nu\) run over spacetime, while Latin indices such as \(i,j\) usually run over the three spatial components only. \lecturetimestamp{01:12:23}}
\end{classroomqa}

The same construction works for two different four-vectors:
\begin{equation}
A^\mu B_\mu
\end{equation}
is a scalar.  One quick proof uses the polarization identity.  Since
\((A+B)^2\) and \((A-B)^2\) are both scalar,
\begin{equation}
(A+B)^2-(A-B)^2=4A^\mu B_\mu
\end{equation}
is scalar as well.  This is the Minkowski version of an ordinary dot product.

\begin{classroomqa}
\audiencequestion{Is \(A^\mu B_\mu\) analogous to the ordinary dot product?}
\lecturerresponse{Exactly, except that the Minkowski metric supplies the relative sign between the time and space components. \lecturetimestamp{01:15:44}}
\end{classroomqa}

\subsection{The gradient is a covector}

Let \(\phi(x)\) be a scalar field, with \(x\) now standing for all four
spacetime coordinates.  Between neighboring events,
\begin{equation}
d\phi
=
\frac{\partial\phi}{\partial x^\mu}dx^\mu
=
\partial_\mu\phi\,dx^\mu.
\label{eq:sr5-dphi}
\end{equation}
The left-hand side is scalar, and \(dx^\mu\) is a four-vector.  Therefore the
coefficient
\begin{equation}
\partial_\mu\phi
=
\left(
\partial_t\phi,
\partial_x\phi,
\partial_y\phi,
\partial_z\phi
\right)
\label{eq:sr5-gradient-lower}
\end{equation}
must transform as a covector.  Raising its index gives
\begin{equation}
\partial^\mu\phi
=
\left(
-\partial_t\phi,
\partial_x\phi,
\partial_y\phi,
\partial_z\phi
\right).
\label{eq:sr5-gradient-upper}
\end{equation}

\begin{figure}[htbp]
\centering
\includegraphics[width=0.90\textwidth]{lecture_05_figure_20.png}
\caption{The differential \(d\phi=\partial_\mu\phi\,dx^\mu\), from which the covector character of the scalar gradient follows.}
\lectureframe{01:20:00}
\end{figure}

\begin{classroomqa}
\audiencequestion{Does the upper index in \(x^\mu\) become a lower index when it appears in the denominator \(\partial/\partial x^\mu\)?}
\lecturerresponse{Yes.  The coordinate differential \(dx^\mu\) transforms contravariantly, so its dual derivative \(\partial_\mu=\partial/\partial x^\mu\) transforms covariantly.  That is what makes \(d\phi=\partial_\mu\phi\,dx^\mu\) a scalar. \lecturetimestamp{01:20:00}}
\end{classroomqa}

\subsection{An invariant scalar-field density}

Contract the two derivative vectors:
\begin{align}
\partial_\mu\phi\,\partial^\mu\phi
&=
-(\partial_t\phi)^2
+(\partial_x\phi)^2
+(\partial_y\phi)^2
+(\partial_z\phi)^2.
\end{align}
The free scalar density can therefore be written compactly as
\begin{equation}
\boxed{
\mathcal L_0
=
-\frac12
\partial_\mu\phi\,\partial^\mu\phi
}.
\label{eq:sr5-compact-density}
\end{equation}
This is identical to the component density in
\eqref{eq:sr5-free-field-action}.

\begin{figure}[htbp]
\centering
\includegraphics[width=0.90\textwidth]{lecture_05_figure_21.png}
\caption{The scalar-field kinetic density condensed to one contracted expression.}
\lectureframe{01:25:10}
\end{figure}

\begin{classroomqa}
\audiencequestion{Is the contracted first-derivative expression the Laplacian?}
\lecturerresponse{No.  The Laplacian contains second spatial derivatives.  This expression is the Lagrangian density, built from products of first derivatives. \lecturetimestamp{01:24:38}}
\end{classroomqa}

The value of a scalar density at a physical event is frame independent.  The
proper Lorentz transformation also has determinant one in magnitude, so the
measure \(d^4x\) is invariant.  Consequently,
\begin{equation}
S=\int d^4x\,\mathcal L
\end{equation}
is invariant.  A frame-independent action yields Euler--Lagrange equations
with the same form in every inertial frame.

\section{The Galilean Limit}

\begin{classroomqa}
\audiencequestion{Was a Newtonian Lagrangian also required to respect a relativity principle?}
\lecturerresponse{Yes, though it was not emphasized earlier.  Newtonian equations are invariant under Galilean transformations, the \(c\to\infty\) limit of Lorentz transformations.  What disappears in that limit is relativity of simultaneity. \lecturetimestamp{01:27:16}}
\end{classroomqa}

To see that limit, restore \(c\) in a boost along \(x\):
\begin{equation}
x'
=
\frac{x-vt}{\sqrt{1-v^2/c^2}},
\qquad
t'
=
\frac{t-vx/c^2}{\sqrt{1-v^2/c^2}}.
\label{eq:sr5-lorentz-c}
\end{equation}
Taking \(c\to\infty\) gives
\begin{equation}
\boxed{
x'=x-vt,
\qquad
t'=t
}.
\label{eq:sr5-galilean}
\end{equation}
These are Galilean transformations.  Newtonian mechanics retains the
relativity principle that the laws have the same form in all inertial frames,
but it also assumes an absolute time shared by every observer.

\begin{figure}[htbp]
\centering
\includegraphics[width=0.90\textwidth]{lecture_05_figure_22.png}
\caption{The Lorentz boost before taking \(c\to\infty\), leaving the Galilean rules \(x'=x-vt\) and \(t'=t\).}
\lectureframe{01:29:40}
\end{figure}

\begin{classroomqa}
\audiencequestion{Do upper- and lower-index vectors transform differently only by a sign change?}
\lecturerresponse{For the fixed Cartesian metric used here, lowering an already known vector changes only its time-component sign.  Their Lorentz transformation laws are nevertheless dual laws, not merely an arbitrary sign convention; the metric is what converts one into the other. \lecturetimestamp{01:31:19}}
\end{classroomqa}

\begin{classroomqa}
\audiencequestion{Is every ordinary-looking function of coordinates a scalar?}
\lecturerresponse{No.  Being one number at each coordinate point is not enough; the number must represent the same physical value at the same event in every frame.  Energy density is a counterexample: boosting a box changes both the measured energy and the measured volume, so observers do not assign the same density. \lecturetimestamp{01:33:23}}
\end{classroomqa}

A single component of four-velocity provides another example.  The full
four-velocity is a vector, but its time component changes under a boost and is
not itself a scalar.

\section{The Klein--Gordon Field}

\subsection{The field analog of a harmonic oscillator}

Return to the scalar density and add the simplest quadratic potential:
\begin{equation}
\boxed{
\mathcal L
=
\frac12
\left[
(\partial_t\phi)^2
-|\boldsymbol{\nabla}\phi|^2
-\mu^2\phi^2
\right]
}.
\label{eq:sr5-kg-density}
\end{equation}
This is the field-theory analog of
\begin{equation}
L_{\mathrm{osc}}
=
\frac12\dot q^2-\frac12\mu^2q^2.
\end{equation}
Every spatial point behaves locally like an oscillator, while the gradient
term couples neighboring points.

\begin{figure}[htbp]
\centering
\includegraphics[width=0.90\textwidth]{lecture_05_figure_23.png}
\caption{The quadratic scalar-field density and its harmonic-oscillator prototype.}
\lectureframe{01:38:20}
\end{figure}

The field Euler--Lagrange equation is
\begin{equation}
\partial_\mu
\left[
\frac{\partial\mathcal L}
{\partial(\partial_\mu\phi)}
\right]
-
\frac{\partial\mathcal L}{\partial\phi}
=0.
\label{eq:sr5-field-el}
\end{equation}
Applying it to \eqref{eq:sr5-kg-density} gives
\begin{equation}
\boxed{
\partial_t^2\phi-\nabla^2\phi+\mu^2\phi=0
}.
\label{eq:sr5-kg}
\end{equation}
This is the Klein--Gordon equation.  With the mostly-plus convention and
\(\Box=\partial_\mu\partial^\mu=-\partial_t^2+\nabla^2\), the same equation is
\((\Box-\mu^2)\phi=0\).

\begin{figure}[htbp]
\centering
\includegraphics[width=0.90\textwidth]{lecture_05_figure_24.png}
\caption{The Klein--Gordon equation completed beneath the density from which it follows.}
\lectureframe{01:40:00}
\end{figure}

Historically, the equation was first explored as a relativistic
single-particle wave equation.  That interpretation encountered difficulties,
especially for a positive probability density.  Its proper modern setting is
relativistic field theory, where it governs free scalar fields and the free
limit of interacting scalar fields.

\subsection{Plane-wave modes}

Because \eqref{eq:sr5-kg} is linear with constant coefficients, seek a plane
wave
\begin{equation}
\phi(x)
=
\exp\!\left[
-i\omega t+i\mathbf{k}\cdot\mathbf{x}
\right].
\label{eq:sr5-plane-wave}
\end{equation}
With
\begin{equation}
k_\mu=(-\omega,k_x,k_y,k_z),
\end{equation}
the exponent is \(ik_\mu x^\mu\).

\begin{figure}[htbp]
\centering
\includegraphics[width=0.90\textwidth]{lecture_05_figure_25.png}
\caption{The plane-wave ansatz written in components and as \(e^{ik_\mu x^\mu}\).}
\lectureframe{01:44:00}
\end{figure}

Differentiation gives
\begin{equation}
\partial_t^2\phi=-\omega^2\phi,
\qquad
\nabla^2\phi=-|\mathbf{k}|^2\phi.
\end{equation}
Substitution into \eqref{eq:sr5-kg} yields
\begin{equation}
\left(
-\omega^2+|\mathbf{k}|^2+\mu^2
\right)\phi=0.
\end{equation}
A nonzero mode therefore requires
\begin{equation}
\boxed{
\omega^2=|\mathbf{k}|^2+\mu^2
},
\qquad
\omega=\pm\sqrt{|\mathbf{k}|^2+\mu^2}.
\label{eq:sr5-dispersion}
\end{equation}

\begin{figure}[htbp]
\centering
\includegraphics[width=0.90\textwidth]{lecture_05_figure_26.png}
\caption{The plane-wave substitution and the resulting relativistic dispersion relation.}
\lectureframe{01:47:00}
\end{figure}

In natural units, \(\mu\) has the role of a mass.  Restoring constants for a
quantum excitation,
\begin{equation}
E=\hbar\omega,
\qquad
\mathbf{p}=\hbar\mathbf{k},
\qquad
\mu=\frac{mc}{\hbar},
\end{equation}
turns \eqref{eq:sr5-dispersion} into
\begin{equation}
E^2=p^2c^2+m^2c^4.
\end{equation}
The positive- and negative-frequency branches are both part of the classical
solution space.  General solutions are superpositions of these modes.

\begin{classroomqa}
\audiencequestion{Where did the Klein--Gordon equation come from?}
\lecturerresponse{It comes directly from the quadratic Lorentz-scalar density through the field Euler--Lagrange equation.  The example was chosen because this simple quadratic density produces a linear field equation. \lecturetimestamp{01:47:59}}
\end{classroomqa}

\begin{classroomqa}
\audiencequestion{What in nature obeys the Klein--Gordon equation?}
\lecturerresponse{A free classical scalar field does.  Higgs-field fluctuations are scalar and reduce to Klein--Gordon form in the free approximation, although the actual Higgs field also has interactions and gauge couplings. \lecturetimestamp{01:50:41}}
\end{classroomqa}

\section{A Scalar Background as a Mass Shift}

The particle action now comes back into view:
\begin{equation}
S_{\mathrm p}
=
-\int d\tau\,
\bigl[m+g\phi(x)\bigr].
\label{eq:sr5-final-particle-action}
\end{equation}
If \(\phi\) has a constant background value \(\phi_0\), the particle moves as
though its mass were
\begin{equation}
\boxed{
m_{\mathrm{eff}}=m+g\phi_0
}.
\label{eq:sr5-effective-mass}
\end{equation}
If the bare term \(m\) vanishes, the scalar background can provide the entire
coefficient multiplying proper time.

\begin{figure}[htbp]
\centering
\includegraphics[width=0.90\textwidth]{lecture_05_figure_27.png}
\caption{The field-dependent particle action above the compact field Euler--Lagrange equation, bringing the two sides of the lecture together.}
\lectureframe{01:51:30}
\end{figure}

This is a deliberately simplified model of Higgs mass generation.  In the
Standard Model, fermion masses arise from Yukawa couplings to the Higgs vacuum
expectation value, schematically \(m_f\propto y_f v\), while the \(W\) and
\(Z\) masses arise through the gauge-covariant Higgs kinetic term.  The toy
worldline coupling captures the central idea--a nonzero scalar background
changes the inertial mass--without reproducing the full electroweak theory.
The reciprocity established earlier remains in force: varying the same
interaction with respect to the field makes the particle a source.  The field
acts on particles, and particles act back on the field.

\subsection{Why complex exponentials are useful}

The field need not be physically complex.  Euler's formula gives
\begin{equation}
e^{i(\omega t-\mathbf{k}\cdot\mathbf{x})}
=
\cos(\omega t-\mathbf{k}\cdot\mathbf{x})
+i\sin(\omega t-\mathbf{k}\cdot\mathbf{x}).
\end{equation}
Since the equation is real and linear, the real and imaginary parts are
separately real solutions.  They are waves displaced by one quarter of a
cycle, and any real linear combination is again a solution.

\begin{classroomqa}
\audiencequestion{When the real part is taken, does the imaginary part somehow add to zero?}
\lecturerresponse{No.  The complex exponential packages two independent real solutions, cosine and sine.  Either part may be retained, or real combinations may be formed.  The exponential is used because differentiation becomes multiplication. \lecturetimestamp{01:53:29}}
\end{classroomqa}

\begin{figure}[htbp]
\centering
\includegraphics[width=0.90\textwidth]{lecture_05_figure_28.png}
\caption{A complex plane wave split into cosine and sine, drawn as two real waves out of phase.}
\lectureframe{01:56:00}
\end{figure}

\begin{classroomqa}
\audiencequestion{Is the mass acquired by the electromagnetic field in a superconductor the same kind of phenomenon?}
\lecturerresponse{It is closely related: in a superconductor the condensate gives the electromagnetic field an effective mass in the medium, producing the Meissner effect.  A full explanation requires the electromagnetic field and is postponed. \lecturetimestamp{01:56:14}}
\end{classroomqa}

\begin{classroomqa}
\audiencequestion{Can the Higgs mechanism be visualized as a Higgs particle flying through space and handing out mass?}
\lecturerresponse{No.  The relevant object is the background value of the Higgs field, not a stream of individual Higgs particles.  In the toy action, the background simply changes the coefficient that functions as the particle's mass. \lecturetimestamp{01:57:31}}
\end{classroomqa}

\begin{classroomqa}
\audiencequestion{Can different couplings to the same scalar background produce different masses?}
\lecturerresponse{Yes.  The shift is proportional to the coupling \(g\).  In the Standard Model, different Yukawa couplings are precisely why different fermions acquire different masses from the same Higgs background. \lecturetimestamp{01:58:04}}
\end{classroomqa}

\begin{classroomqa}
\audiencequestion{Is \(\phi\) itself the mass term?}
\lecturerresponse{In this toy model, \(g\phi_0\) functions in every mechanical respect as a contribution to the mass coefficient.  The coupling and the background value must both be included. \lecturetimestamp{01:58:42}}
\end{classroomqa}

\begin{classroomqa}
\audiencequestion{If the Higgs background vanished, would every particle become massless?}
\lecturerresponse{No.  Electrons, quarks, and the weak gauge bosons would lose their ordinary electroweak mass contributions, with neutrino details depending on the mechanism that gives neutrinos mass.  The scalar sector itself need not become massless merely because its background value vanishes.  Most of the proton and neutron mass comes instead from QCD binding and confinement, so it is not generated directly by the Higgs field. \lecturetimestamp{01:58:51}}
\end{classroomqa}

The atomic consequence is nevertheless drastic.  The Bohr scale behaves as
\begin{equation}
a_0\sim\frac{1}{\alpha m_e}
\end{equation}
in natural units.  Sending the electron mass toward zero makes the ordinary
atomic length scale diverge; ordinary localized atoms would cease to exist.
The proton and neutron would be much less directly affected because their
mass is predominantly dynamical QCD energy.

The pieces now form one continuous argument.  Lorentz
geometry controls observed wavelengths.  One mixed term in one action makes
two dynamical systems influence each other.  A scalar field becomes both a
potential for a particle and a field sourced by that particle.  The metric
and index notation make Lorentz invariance visible, and the simplest
quadratic invariant density produces Klein--Gordon waves.  Finally, a
constant scalar background turns the same coupling into a mass shift.  The
method throughout is the one that will persist: construct an invariant
action, vary every dynamical object, and read the mutual physics from the
resulting equations.
